% =============================================================================
% Capítulo 5 — Avaliação Quantitativa
% =============================================================================
\chapter{Avaliação Quantitativa}
\label{cap:avaliacao}

Este capítulo descreve o cenário, as condições e o protocolo adotados para a
avaliação quantitativa comparativa entre o chefe controlado por FSM e o chefe
controlado pelo modelo PPO treinado. O objetivo da avaliação é verificar se o
pipeline de \textit{Reward Shaping} por proveniência produziu um agente com
comportamento mensuravelmente distinto do \textit{baseline} determinístico,
utilizando métricas objetivas coletadas em sessões controladas. Os resultados
numéricos da avaliação são apresentados no \autoref{cap:resultados}.

% -----------------------------------------------------------------------------
\section{Ambiente de Avaliação}
\label{sec:ambiente-avaliacao}

A avaliação foi conduzida em duas cenas do Unity, cada uma configurada para
uma das condições experimentais. Ambas as cenas compartilham a mesma arena de
combate, os mesmos parâmetros de dano e vida, e o mesmo personagem jogável
(Elian), diferindo exclusivamente no componente que controla o chefe.

A \autoref{tab:cenas-oficiais} apresenta as cenas oficiais utilizadas.

\begin{table}[htb]
    \centering
    \caption{Cenas oficiais de avaliação e suas configurações.}
    \label{tab:cenas-oficiais}
    \begin{tabular}{l p{5cm} p{5.5cm}}
        \toprule
        \textbf{Condição} & \textbf{Cena Unity}
                          & \textbf{Configuração} \\
        \midrule
        FSM & \texttt{ArenaChefe.unity}
            & Chefe controlado por \texttt{BossFsmController}. Exporta CSV
              de métricas e JSON de proveniência. \\
        \addlinespace
        PPO & \texttt{BossArena\_PPO.unity}
            & Chefe controlado por \texttt{BossAgent} em modo
              \texttt{InferenceOnly}. Exporta CSV de métricas. \\
        \bottomrule
    \end{tabular}
    \fonte{Elaborado pelo autor.}
\end{table}

\subsection{Preparação das Cenas}
\label{subsec:preparacao-cenas}

Ambas as cenas foram preparadas para serem jogáveis pelo desenvolvedor em
condições idênticas. Essa preparação envolveu ajustes específicos para
garantir a equivalência entre as condições experimentais:

\begin{itemize}
    \item \textbf{Personagem controlável}: nas duas cenas, Elian é controlado
          pelo desenvolvedor via \textit{Input System}, com as mesmas
          mecânicas de movimento, pulo, ataque e \textit{dash};
    \item \textbf{Câmera}: o componente \texttt{HorizontalCameraFollow} foi
          adicionado para que a câmera acompanhe o personagem apenas no eixo
          horizontal, melhorando a jogabilidade sem alterar nenhum aspecto da
          IA, do sistema de recompensa, do modelo, das \textit{hitboxes} ou
          das métricas;
    \item \textbf{Parâmetros equalizados}: os valores de vida do jogador, dano
          do chefe, alcance de ataque e \textit{cooldown} de dano são
          idênticos nas duas cenas, eliminando vieses de configuração.
\end{itemize}

A cena PPO passou por uma adaptação adicional: durante o desenvolvimento, ela
era orientada ao treinamento e utilizava um \texttt{PlayerDummyAI} como
oponente simulado. Para a fase de avaliação, a dependência do
\texttt{PlayerDummyAI} foi removida e a cena foi reconfigurada para receber
entrada humana, garantindo que as sessões contra o chefe PPO seguissem o mesmo
fluxo operacional das sessões contra o chefe FSM.

\subsection{Exportação de Proveniência na Avaliação}
\label{subsec:proveniencia-avaliacao}

Uma distinção operacional entre as duas cenas é a exportação de grafos de
proveniência em JSON. Na cena FSM, a exportação ao término da sessão
(\texttt{exportOnSessionEnd}) permaneceu habilitada, gerando um arquivo JSON
por sessão. Na cena PPO, essa exportação foi desabilitada por decisão de
escopo: a métrica principal da avaliação comparativa é o CSV quantitativo, e
os JSONs da cena FSM já constituem evidência suficiente do funcionamento do
mecanismo de proveniência.

Essa decisão não afeta a validade da comparação quantitativa, pois ambas as
cenas exportam o mesmo CSV de métricas com estrutura de colunas idêntica.

% -----------------------------------------------------------------------------
\section{Condições de Vitória e Derrota}
\label{sec:condicoes-vitoria-derrota}

A definição de condições reais de vitória e derrota é um requisito
metodológico fundamental para a validade da avaliação. Uma comparação entre
chefes FSM e PPO na qual o jogador não pudesse ser derrotado produziria dados
enviesados: o campo \texttt{player\_damage\_taken} seria artificial, o
resultado de todas as sessões tenderia invariavelmente a vitória, e a
capacidade ofensiva do chefe --- que é precisamente a variável de interesse
--- não seria adequadamente capturada pelas métricas.

\subsection{Adição do Sistema de Vida Real}
\label{subsec:adicao-vida-real}

Para resolver essa limitação, o componente \texttt{PlayerHealth} foi integrado
ao personagem Elian antes do início da coleta oficial. Esse componente
implementa um sistema de pontos de vida com valor máximo de 100, dano
recebido cumulativo e emissão de evento de morte quando os pontos de vida
atingem zero. A morte do jogador é detectada pelo \texttt{ArenaManager}, que
encerra a sessão com resultado \texttt{defeat}.

Complementarmente, o componente \texttt{BossAttackDamage} foi implementado
para aplicar dano real ao jogador quando o chefe executa um ataque dentro do
alcance. Esse componente é compartilhado pelas duas condições experimentais
(FSM e PPO), garantindo que o sistema de dano opere de forma idêntica
independentemente do controlador do chefe.

\subsection{Parâmetros de Dano}
\label{subsec:parametros-dano}

A \autoref{tab:parametros-dano} apresenta os parâmetros de dano utilizados na
avaliação.

\begin{table}[htb]
    \centering
    \caption{Parâmetros de dano do chefe utilizados na avaliação.}
    \label{tab:parametros-dano}
    \begin{tabular}{l r l}
        \toprule
        \textbf{Parâmetro}                 & \textbf{Valor} & \textbf{Unidade} \\
        \midrule
        \texttt{PlayerHealth.maxHealth}    & 100   & pontos de vida \\
        \texttt{damageAmount}              & 10    & pontos por ataque \\
        \texttt{attackRange}               & 1,25  & unidades \\
        \texttt{damageCooldown}            & 0,8   & segundos \\
        \bottomrule
    \end{tabular}
    \fonte{Elaborado pelo autor.}
\end{table}

Todos os parâmetros foram mantidos idênticos nas duas condições para
eliminar vieses de configuração. A igualdade do \texttt{damageAmount} e do
\texttt{attackRange} garante que cada ataque do chefe, seja FSM ou PPO,
produza exatamente o mesmo efeito quando conectado.

\subsection{O Papel do Cooldown de Dano}
\label{subsec:cooldown-dano}

O parâmetro \texttt{damageCooldown} merece atenção especial. Esse
\textit{cooldown} define o intervalo mínimo entre duas aplicações consecutivas
de dano pelo chefe, independentemente da frequência com que o chefe executa
ações de ataque.

No contexto da condição FSM, o \textit{cooldown} de dano é pouco perceptível,
pois o \texttt{attackCooldown} interno do FSM (também de 0,8 segundos) já
limita naturalmente a frequência de ataques. Na condição PPO, contudo, o
\textit{cooldown} de dano é crucial: o agente treinado apresenta alta
frequência de ações de ataque, executando a ação \textit{Attack Melee} em
uma fração significativa dos passos de decisão. Sem o \textit{cooldown} de
dano, essa frequência elevada resultaria em múltiplas aplicações de dano
por segundo, eliminando o jogador quase instantaneamente e inviabilizando
sessões jogáveis.

O valor de 0,8 segundos foi escolhido para alinhar-se ao \textit{cooldown}
natural do FSM, garantindo que a taxa máxima de dano efetivo seja a mesma nas
duas condições. Dessa forma, a diferença na frequência de ações de ataque
entre FSM e PPO manifesta-se na contagem de ações registrada (\texttt{boss\_attack\_count}),
mas não na taxa de dano efetivo aplicado ao jogador.

\subsection{Integração com os Sistemas Existentes}
\label{subsec:integracao-sistemas}

A adição do sistema de vida real exigiu integração com três sistemas
existentes:

\begin{enumerate}
    \item \textbf{ArenaManager}: a morte do jogador (evento emitido pelo
          \texttt{PlayerHealth}) passou a ser detectada como condição de
          encerramento da sessão, com resultado \texttt{defeat};
    \item \textbf{GameMetrics}: o campo \texttt{player\_damage\_taken} passou
          a ser alimentado pelo dano real aplicado pelo
          \texttt{BossAttackDamage}, em vez de valores artificiais;
    \item \textbf{ProvenanceLogger}: o evento \texttt{PlayerDamageTaken} foi
          adicionado ao conjunto de eventos capturados, registrando cada
          instância de dano real no grafo de proveniência da sessão.
\end{enumerate}

Essas integrações seguem o mesmo padrão de comunicação por eventos descrito
no \autoref{cap:implementacao}: cada sistema inscreve-se nos eventos emitidos
pelo \texttt{PlayerHealth} sem modificar o componente emissor, preservando o
baixo acoplamento da arquitetura.

% -----------------------------------------------------------------------------
\section{Protocolo de Coleta}
\label{sec:protocolo-coleta}

O protocolo de coleta foi projetado para maximizar a comparabilidade entre as
condições experimentais dentro das restrições de escopo do trabalho.

\subsection{Procedimento}
\label{subsec:procedimento}

A coleta oficial consistiu em 20 sessões jogadas manualmente pelo
desenvolvedor, divididas em dois blocos sequenciais:

\begin{enumerate}
    \item \textbf{Bloco FSM}: 10 sessões consecutivas na cena
          \texttt{ArenaChefe.unity}, contra o chefe controlado por FSM;
    \item \textbf{Bloco PPO}: 10 sessões consecutivas na cena
          \texttt{BossArena\_PPO.unity}, contra o chefe controlado pelo
          modelo PPO em modo \texttt{InferenceOnly}.
\end{enumerate}

Em ambos os blocos, o desenvolvedor jogou normalmente como Elian, utilizando
as mecânicas disponíveis (movimento, pulo, ataque e \textit{dash}) sem
restrições ou instruções comportamentais específicas. Cada sessão encerrou-se
naturalmente por vitória (morte do chefe) ou derrota (morte do jogador), e
as métricas foram exportadas automaticamente ao término da sessão.

A \autoref{tab:amostra} sintetiza a composição da amostra oficial.

\begin{table}[htb]
    \centering
    \caption{Composição da amostra oficial da avaliação.}
    \label{tab:amostra}
    \begin{tabular}{l r}
        \toprule
        \textbf{Condição} & \textbf{Sessões} \\
        \midrule
        FSM                & 10 \\
        PPO                & 10 \\
        \midrule
        \textbf{Total}     & \textbf{20} \\
        \bottomrule
    \end{tabular}
    \fonte{Elaborado pelo autor.}
\end{table}

\subsection{Coleta Automatizada de Métricas}
\label{subsec:coleta-automatizada}

As métricas de cada sessão foram registradas automaticamente pelo sistema
\texttt{GameMetrics} e exportadas pelo \texttt{GameMetricsExporter} para o
arquivo CSV \texttt{boss\_evaluation\_metrics.csv}. Esse processo não requer
intervenção manual do jogador: os contadores são inicializados
automaticamente ao início da sessão e a linha CSV é gravada ao término,
garantindo a lisura dos dados coletados.

O cabeçalho do CSV contém 18 campos, abrangendo identificação da sessão
(\texttt{session\_id}, \texttt{boss\_type}), resultado
(\texttt{result}), temporização (\texttt{start\_time\_seconds},
\texttt{end\_time\_seconds}, \texttt{duration\_seconds}), métricas de
combate (\texttt{boss\_damage\_taken}, \texttt{player\_damage\_taken}) e
contagens detalhadas de ações do jogador e do chefe. Para sessões contra o
chefe PPO, seis campos adicionais registram a distribuição das ações
discretas do agente (\texttt{ppo\_idle\_count} a \texttt{ppo\_dash\_count}).

Três cuidados de engenharia garantem a confiabilidade do CSV exportado:

\begin{enumerate}
    \item O cabeçalho é escrito uma única vez, na criação do arquivo,
          assegurando estrutura de colunas consistente;
    \item Valores numéricos são formatados com
          \texttt{CultureInfo.InvariantCulture}, de modo que o separador
          decimal seja sempre o ponto, independentemente da configuração
          regional do sistema;
    \item O caminho do arquivo é resolvido relativamente à raiz do
          repositório, e o diretório de destino é criado automaticamente
          caso não exista.
\end{enumerate}

\subsection{Controle de Linhas Extras}
\label{subsec:controle-linhas}

Após a coleta oficial das 20 sessões, sessões adicionais foram realizadas
para fins de gravação de evidências visuais (vídeos demonstrativos do
comportamento dos chefes FSM e PPO). Essas sessões adicionais geraram linhas
extras no CSV, que não fazem parte da amostra oficial.

Para evitar contaminação dos dados, a análise apresentada no
\autoref{cap:resultados} considera exclusivamente as primeiras 20 linhas do
CSV, correspondentes às 10 sessões FSM e 10 sessões PPO da coleta oficial.
Qualquer linha posterior é descartada da análise.

\subsection{Considerações sobre Validade}
\label{subsec:validade}

O protocolo adotado apresenta características que favorecem a validade interna
da comparação:

\begin{itemize}
    \item \textbf{Condições equalizadas}: parâmetros de dano, vida e alcance
          são idênticos nas duas cenas, garantindo que diferenças nas
          métricas reflitam exclusivamente o comportamento do chefe;
    \item \textbf{Coleta automatizada}: a exportação de métricas não depende
          de registro manual, eliminando erros de transcrição;
    \item \textbf{Modelo congelado}: o chefe PPO opera em inferência pura, sem
          ajustes \textit{online}, garantindo reprodutibilidade do
          comportamento;
    \item \textbf{Mesmo jogador}: todas as sessões foram conduzidas pelo
          mesmo operador, eliminando variância interpessoal.
\end{itemize}

As limitações do protocolo --- em particular, a amostra reduzida, a avaliação
conduzida pelo próprio desenvolvedor e a ausência de significância estatística
formal --- são reconhecidas e discutidas no \autoref{cap:discussao}.
